\section{Kết quả xây dựng}

Kết quả xây dựng của hệ thống được lưu tại GitHub repository: \url{https://github.com/fabyanbui/feature-flag-platform}. Các minh chứng chi tiết về Backend API, Frontend Dashboard, Demo Application và SDK sẽ được trình bày chi tiết trong phần demo.

\subsection{Backend API}

\begin{figure}[H]
    \centering
    \includegraphics[width=0.95\textwidth]{\detokenize{Images/screenshot/Kết quả xây dựng Backend API.png}}
    \caption{Minh chứng kết quả xây dựng Backend API}
    \label{fig:backend-api-result}
\end{figure}

\subsection{Frontend Dashboard}

\begin{figure}[H]
    \centering
    \includegraphics[width=0.95\textwidth]{\detokenize{Images/screenshot/Kết quả xây dựng Frontend Dashboard.png}}
    \caption{Minh chứng kết quả xây dựng Frontend Dashboard}
    \label{fig:frontend-dashboard-result}
\end{figure}

\subsection{Demo Application}

Demo Application đã được mở rộng từ ứng dụng frontend thuần thành storefront có kèm demo backend API nhẹ. Mỗi feature flag trong demo có một endpoint GET tương ứng dưới \path|/api/demo/features/*|; endpoint này dùng \texttt{@ffp/js-sdk} ở phía server để gọi Evaluation API trước khi trả dữ liệu tính năng. Nếu flag tắt, bị archive, bị group/flag kill switch hoặc SDK không xác minh được trạng thái bật, demo backend trả \texttt{403 FEATURE\_DISABLED}.

Một số thao tác có sẵn trong giao diện storefront được nối với endpoint backend tương ứng với feature flag, không cần thêm panel minh họa riêng. Nhờ đó phần demo thể hiện đồng thời hai lớp tích hợp: frontend ẩn/hiện UI theo flag và backend API vẫn tự guard khi bị gọi trực tiếp.

\begin{figure}[H]
    \centering
    \includegraphics[width=0.95\textwidth]{\detokenize{Images/screenshot/Kết quả xây dựng Demo Application.png}}
    \caption{Minh chứng kết quả xây dựng Demo Application}
    \label{fig:demo-application-result}
\end{figure}
